% MT016 / explanation — canonical editable content.
设\(\phi(x)\)把原输入变成一组新特征。内积把两组特征的对应分量相乘再相加；核函数直接返回这个结果：\(k(x,z)=\phi(x)^T\phi(z)\)。我们可以不显式生成很长的\(\phi(x)\)，这就是核技巧。


把每对训练样本的核值排成一张表，得到核矩阵（Gram matrix）：第\(i\)行、第\(j\)列是\(K_{ij}=k(x_i,x_j)\)。例如两条样本对应
\[
K=\begin{pmatrix}k(x_1,x_1)&k(x_1,x_2)\\k(x_2,x_1)&k(x_2,x_2)\end{pmatrix}.
\]


矩阵乘一个向量，就是分别用每一行与该向量求内积；下面C的反例会用到这个运算。



\begin{itemize}

\item \textbf{A正确。} 核SVM对应在映射后的特征空间学习线性分隔，在原空间看起来可能是曲线。

\item \textbf{B错误。} 一般非线性核预测为\(f(x)=\sum_{i\in SV}\alpha_i y_i k(x_i,x)+b\)。新点要与支持向量\(x_i\)算核值，因此通常仍需保存这些训练点及系数；其余训练点可以不参与预测。

\item \textbf{C错误。} 核不只是一个数值为正的相似度函数，它必须对任意有限样本形成对称半正定Gram矩阵。半正定要求：用任意实数系数\(v\)加权组合这些样本，都有\(v^TKv\ge0\)；对内积核，这个量就是组合后向量的平方长度。反例\(K=\begin{pmatrix}1&2\\2&1\end{pmatrix}\)虽然每项都正，取\(v=(1,-1)^T\)却有\(Kv=(-1,1)^T\)，于是\(v^TKv=-1-1=-2\)。平方长度不可能为负，所以它不是合法的内积矩阵。

\item \textbf{D正确。} 字符串或集合也能定义合法核，只要满足内积结构的条件；不要求显式写出有限维向量。

\item \textbf{E正确，理解为可能的节省。} 若映射包含数量庞大的特征，直接算核值可省去它们的构造与存储。不过样本很多时，\(N\times N\)核矩阵本身仍可能很大。

\end{itemize}


合法性与效率是两个问题：满足核条件不意味着一定省内存，输出全正也不意味着满足核条件。半正定的具体检查见\StudyRef{Answers}{mt055}{MT055}。


\textbf{出处：} 2021A期中 Q3，原卷第2页（5分）


\textbf{答案依据：} 2021A期中 Q3 随卷／配套参考答案；条件与解释经复核。
